\documentclass[14pt]{extarticle}

% Esto es para poder escribir acentos directamente:
\usepackage[utf8]{inputenc}
% Esto es para que el LaTeX sepa que el texto est en espaol:
\usepackage[spanish, activeacute]{babel}

% Paquetes de la AMS:
\usepackage{amsmath, amsthm, amsfonts}

\usepackage{extsizes}

\usepackage{pdflscape}
\usepackage{fontawesome}

\usepackage[round]{natbib}

\usepackage{graphicx}
\graphicspath{ {./images/} }

\usepackage[margin=1in]{geometry}

\usepackage{listings}
\usepackage{natbib}
\usepackage{url}
\usepackage{hyperref} % For hyperlinks in the PDF
\usepackage{multirow}
\usepackage{xifthen}
\usepackage{tabularx}

% Teoremas
%--------------------------------------------------------------------------
\newtheorem{thm}{Teorema}[section]
\newtheorem{cor}[thm]{Corolario}
\newtheorem{lem}[thm]{Lema}
\newtheorem{prop}[thm]{Proposicin}
\theoremstyle{definition}
\newtheorem{defn}[thm]{Definicin}
\theoremstyle{remark}
\newtheorem{rem}[thm]{Observacin}

% Atajos.
% Se pueden definir comandos nuevos para acortar cosas que se usan
% frecuentemente. Como ejemplo, aqu se definen la R y la Z dobles que
% suelen representar a los conjuntos de nmeros reales y enteros.
%--------------------------------------------------------------------------

\def\RR{\mathbb{R}}
\def\ZZ{\mathbb{Z}}

% De la misma forma se pueden definir comandos con argumentos. Por
% ejemplo, aqu definimos un comando para escribir el valor absoluto
% de algo ms fcilmente.
%--------------------------------------------------------------------------
\newcommand{\abs}[1]{\left\vert#1\right\vert}

\newcommand{\newtableitem}[1] {
	\multicolumn{4}{|l|}{· #1} \\
}

\newcommand{\newtableline}[1] { %55 caracteres
	\multicolumn{4}{|l|}{#1} \\
}

\newcommand{\userstory}[9]{
	\begin{center}
	\resizebox{15cm}{!} {
	\begin{tabular}{|c|c|c|c|}
    	\hline 
		\multicolumn{4}{|c|}{\textbf{#1}} \\
		\multicolumn{4}{|c|}{} \\
		\hline	
		\textbf{Identificador:} #2 & \textbf{Prioridad:} #3 & \textbf{Iteración:} #4 & \textbf{Puntos de historia:} #5 \\ 
		\hline 
		\multicolumn{4}{|l|}{\textbf{Descripción:}} \\
		#6
		\hline 
		\multicolumn{4}{|l|}{\textbf{Tareas:}} \\
		#7
		\hline 
		\multicolumn{4}{|l|}{\textbf{Pruebas de aceptación:}} \\
		#8
		\hline 
		\multicolumn{4}{|l|}{\textbf{Observaciones:}} \\
		#9
		\hline 
	\end{tabular}
	}
	\end{center}
}

\newcommand{\sprinttask}[2] {
	\multicolumn{1}{|>{\hsize=\dimexpr1\hsize+1\tabcolsep+\arrayrulewidth\relax}X|}{#1} & \multicolumn{3}{|>{\hsize=\dimexpr3\hsize+3\tabcolsep+\arrayrulewidth\relax}X|}{#2}\\
}

\newcommand{\sprint}[4]{
	\begin{center}
	\begin{tabularx}{\textwidth}{|X|X|X|X|}
    	\hline 
		\multicolumn{2}{|>{\hsize=\dimexpr2\hsize+2\tabcolsep+\arrayrulewidth\relax}X|}{\textbf{Sprint:} #1} &
		\multicolumn{2}{|>{\hsize=\dimexpr2\hsize+2\tabcolsep+\arrayrulewidth\relax}X|}{\textbf{Inicio:} #2} \\
		\hline 
		Versión & \multicolumn{3}{|>{\hsize=\dimexpr3\hsize+3\tabcolsep+\arrayrulewidth\relax}X|}{#4} \\
		\hline 
		Identificador & \multicolumn{3}{|>{\hsize=\dimexpr3\hsize+3\tabcolsep+\arrayrulewidth\relax}X|}{Historia de usuario} \\
		\hline
		#3
		\hline 
	\end{tabularx}
	\end{center}
}

% Operadores.
% Los operadores nuevos deben definirse como tales para que aparezcan
% correctamente. Como ejemplo definimos en jacobiano:
%--------------------------------------------------------------------------
\DeclareMathOperator{\Jac}{Jac}

%--------------------------------------------------------------------------
%\title{Trabajo Fin de Máster \\  TweetSC: Corrector de texto para twitter.}
%\author{Javier Moreno}
\begin{document}
%\maketitle
\begin{titlepage}

\begin{center}
\vspace*{-1in}
\begin{figure}[htb]
\begin{center}
\includegraphics[scale=0.8]{logoUPM.jpg}
\end{center}
\end{figure}
ESCUELA TÉCNICA SUPERIOR DE INGENIEROS INFORMÁTICOS\\
\vspace*{0.15in}
MÁSTER UNIVERSITARIO EN INTELIGENCIA ARTIFICIAL \\
\vspace*{0.6in}
\begin{large}
TRABAJO DE FIN DE MÁSTER:\\
\end{large}
\vspace*{0.2in}
\begin{Large}
\textbf{TWEETSC: CORRECTOR DE TEXTO PARA TWITTER} \\
\end{Large}
\vspace*{0.3in}
\begin{large}
JAVIER MORENO VEGA\\
\end{large}
\vspace*{0.3in}
\rule{80mm}{0.1mm}\\
\vspace*{0.1in}
\begin{large}
TUTOR DE PROYECTO: \\
OSCAR CORCHO GARCÍA \\
\vspace*{0.3in}
CO-TUTOR DE PROYECTO: \\
VÍCTOR RODRÍGUEZ DONCEL \\
\end{large}
\vspace*{2in}
\url{http://tweetsc.github.io}
\vspace*{0.2in}\\
\today
\end{center}
\end{titlepage}
\newpage
\tableofcontents % índice de contenidos
\cleardoublepage
\section{Introducción}\label{sec:introduccion}
\subsection{Motivación}\label{sec:motivacion}
Los nuevos sistemas de comunicación como la mensajería instantánea, chats, redes sociales han generado un uso diferente de los idiomas en estos ámbitos, llamado lenguaje tipo chat \citep{forsyth:2007}. Una de estas redes sociales y en la que este trabajo va a centrarse es Twitter. En esta red social predomina el Uso de emoticonos, repetición de vocales o eliminación de las mismas, uso abusivo de mayúsculas o asusencia, siglas de expresiones populares; lo que dificulta el análisis de los textos. Las ventajas que ofrece esta red social para investigar sobre ella son la cantidad de datos en tiempo real y su fácil acceso.\\\\Uno de los principales problemas a la hora de analizar textos procedentes de las redes sociales son los errores gramaticales que suelen contener, así como la presencia de elementos propios de este tipo de foros que requieren de un procesamiento especial (i.e. hashtags, formas de mencionar a otros usuarios o emoticonos y expresiones habituales en las redes). Además, la limitación en el número de caracteres existente en Twitter la convierte en un caso singular dentro de las redes sociales, ya que los usuarios tienden a adaptar su forma de escribir a dicha limitación, omitiendo palabras y creando abreviaturas que dificultan el uso de herramientas genéricas de procesamiento del lenguaje, especialmente a la hora de realizar tareas como el Análisis de Sentimientos.\\\\Los usuarios en twitter tienden a cometer errores tipográficos, abreviaciones, sustituciones fonéticas y estructuras no gramaticales en los mensajes cortos de texto, causando problemas en las herramientas de análisis. Esto es lo que se consideran palabras mal formadas y la detección de las palabras mal formadas es difícil debido al contexto ruidoso. El objetivo es normalizar estas palabra mal formadas.\\A parte de un uso puramente de investigación, este tipo de trabajo también es beneficioso para un estudio de marcas o personas y sobre lo que las persones opinan sobre ello en las redes social, ya que sin el proceso de normalización y análisis de sentimientos estariamos ante millones de datos que costarían mucho trabajo analizar de una forma automática.

\subsection{Objetivos}\label{sec:objetivos}
El objetivo principal de este trabajo es la creación de un corrector que "normalice" tweets en español.\\Se ha construido un módulo software cuya funcionalidad es:
\begin{itemize}
\item Corregir tweets (palabras y gramaticalmente)
\item Procesar emoticonos deduciendo su significado en el contexto de análisis de sentimientos
\item Expandir hashtags
\item Procesar las conversaciones de los tweets, así como las URLs o imágenes para añadir contexto
\end{itemize}
Todo este módulo software estará disponible a través de una aplicación web.

\subsection{Resumen del documento}\label{sec:resumen}
TODO

\section{Estado del arte}\label{sec:estadodelarte}
subsecciones

Inicio
normalizacion
adaptación de herramientas
Aplicaciones empresariales


En la actualidad, la normalización lingüística de tweets \citep{baldwin:2011} supone un campo de gran interés y en donde la mayoría de trabajos se han realizado sobre textos en inglés y pocos en español. Además no hay ningún trabajo en donde se incluya, dentro de la normalización de tuits, el estudio de los hashtags o etiquetas y los emoticonos, y su contexto. 
Una introducción al tema de normalización de tuits es el artículo \citep{eisenstein:2013}, donde se revisa el estado del arte en NLP sobre variantes SMS y tweets, y cómo la comunidad científica ha respondido por dos caminos: normalización y adaptación de herramientas.\\\\
El modelo de canal ruidoso (Shannon, 1948) ha sido tradicionalmente la primera aproximación a la normalización de textos. Supone que el texto mal formado es T y su forma normalizada es S, por lo que hay que encontrar: arg max P(S|T), calculando arg max P(T|S) P(S), P(S) es el modelo del lenguaje y P(T|S) es el modelo de error. Brill y Moore (2000) caracterizan el modelo de error calculando el producto de operaciones de probabilidad en partes de cadenas de caracteres. Toutanova y Moore (2002) mejoraron el modelo incorporando información de la pronunciación. Choudhury et al. (2007) modela el proceso de generación de texto a nivel de palabra para mensajes SMS considerando las abreviaturas grafémicas/fonéticas y los errores tipográficos involuntarios como transiciones de estado ocultas del modelo de Markov (HMM) y emisiones, respectivamente. Cook y Stevenson (2009) expandieron el modelo de error introduciendo inferencias de diferentes procesos de formación erróneos, de acuerdo con la distribución de errores muestreada. Mientras el modelo de canal ruidoso es apropiado para normalización de textos, es difícil aproximar la normalización con exactitud, además estos métodos ignoran el contexto alrededor del OOV, el cual ayuda a resolver ambigüedades. La traducción automática estadística (SMT) se ha propuesto como un medio de normalización de texto sensible al contexto, al tratar el texto mal formado como el idioma de origen, y la forma estándar como el idioma de destino. Por ejemplo Aw et al. (2006). Normalización de textos como un problema de reconocimiento de voz Kobus et al. (2008). Beufort et al. (2002) métodos de estado finitos combinando las ventajas de SMS y el modelo de canal ruidoso. Kaufmann y Kalita (2010) usan un enfoque de traducción automática con un preprocesador para la normalización sintáctica (en lugar de léxica). El problema de estos trabajos anteriores es que requieren datos de entrenamiento anotados a gran escala, lo que limita su adaptabilidad a nuevos dominios o idiomas, mientras que el trabajo \citep{baldwin:2011} no. Este trabajo es una buena referencia en el campo de la normalización de tuits en inglés. En donde para detectar palabras fuera de diccionario (OOV) utilizan GNU aspell, y los usuarios (@usuario), los hashtags y las URLs son excluidas de la normalización. La normalización tiene relación con los correctores de texto \citep{peterson:1980} pero difiere en que las palabras mal formadas en los mensajes de texto suelen ser intencionadas, para ahorrar caracteres, como identidad social, o debido a la convención en este subgénero de texto. La detección de las palabras mal formadas es difícil debido al contexto ruidoso. El objetivo es normalizar estas palabra mal formadas, además muchas palabras mal formadas son ambiguas y requieren el contexto para poder normalizarlas.\\\\En vez de adaptar el texto a herramientas de análisis, adaptar las herramientas de análisis al texto. Reconocimiento de voz (Gimpel et al., 2011; Owoputi et al., 2013), reconocimiento de entidades (Finin et al., 2010; Ritter et al., 2011; Liu et al., 2011), análisis gramatical (Foster et al., 2011), modelización de diálogos (Ritter et al., 2010) y resumen (Sharifi et al., 2010). El trabajo \citep{liu:2011} sobre NER (reconocimiento de entidades) que replantea el tema de reconocimiento de entidades nombradas en corpus de tuits. Combina un clasificador KNN con CRF (Conditional Random Fields).
\\\\Una introducción a la normalización de tuits en español es \citep{alegria:2013}\citep{alegria:2015}. Utiliza la herramienta Freeling \citep{freeling} para detectar palabras OOV. Uno de los sistemas de normalización de tuits en español, que participó en Tweet-Norm 2013 \citep{alegria:2013}, es \citep{ruizcuadros:2013}\citep{vicomtech}, que usa reglas de preproceso, un modelo de distancias de edición adecuado al dominio y modelos de lengua para seleccionar candidatos de corrección según el contexto. El sistema obtuvo resultados superiores a la media en la tarea \citep{alegria:2013}\citep{tweetnorm}. Una mejora a este trabajo por los mismos autores es \citep{ruizcuadros:2014}. El trabajo que mejores resultados obtuvo en Tweet-Norm 2013 fue RAE \citep{porta:2013}, consiste en transductores que se aplican a tokens OOV. Los transductores implementan modelos lingüísticos de variación que generan conjuntos de candidatos de acuerdo con un léxico. Se usa un modelo de lenguaje estadístico para obtener la secuencia de palabras más probable. En el trabajo \citep{cotelocruz:2015} hace uso de una combinación de varios “módulos expertos” independientes, cada uno especializado en una tarea concreta de la normalización de tuits, en lugar de centrarse en una sola técnica. En este trabajo además realiza un estado del arte actual de la normalización de tuits y en concreto para el idioma español. Otros trabajos sobre normalización en español son (Gomez-Hidalgo, Caurcel-Díaz, y del Rio, 2013) (Mosquera, Lloret, y Moreda, 2012) (Oliva et al., 2011) principalmente sobre mensajes SMS, pero que no abordan la normalización de tuits en su conjunto.
\\\\
Un campo muy relacionado con la normalización de tuits es el análisis de sentimientos y un trabajo que realiza un estudio sobre técnicas de análisis de sentimientos de tuits en español es \citep{fernandezanta:2013}. El trabajo \citep{gamallo:2013} se centra en una técnica Naive-Bayes para el análisis de sentimientos en tuits en español.


\section{Análisis y diseño}\label{sec:analisisydiseno}
\subsection{Metodología de desarrollo}\label{sec:metodologiadedesarrollo}
\subsection{Análisis de requisitos}\label{sec:analisisderequisitos}
\subsection{Solución propuesta}\label{sec:solucionpropuesta}
\section{Implementación}\label{sec:implementacion}
\subsection{Introducción}\label{sec:introduccion}
\subsection{Readme Github}\label{sec:readmegithub}
\subsection{Javadoc}\label{sec:javadoc}
\section{Evaluación}\label{sec:evaluacion}
\subsection{Metodología}\label{sec:metodologia}
\subsection{Corpus}\label{sec:corpus}
\subsection{Gold Standard}\label{sec:goalstandard}
\subsection{Experimentos}\label{sec:experimentos}
\section{Apéndices}\label{sec:apendices}
\subsection{Apéndice A: Bibliografía}\label{sec:bibliografia}
\bibliographystyle{plainnat}
\bibliography{References}

\newpage
\subsection{Apéndice B: Glosario de Términos}\label{sec:glosariodeterminos}
\end{document}